% Ensemble-only manuscript replacement blocks for template_revision.tex
% These blocks are written to replace clustering-focused wording with the
% final ensemble-only CDC method. They are intended to be pasted into the
% manuscript, not cited as a supplementary file.

% -------------------------------------------------------------------------
% Replace line 155 paragraph
% -------------------------------------------------------------------------
Here we directly address these methodological challenges by implementing three algorithmic improvements to our earlier approach: (i) preserving significant local optima from each Monte Carlo realisation (local maxima in goodness, equivalently local minima in K--S distance) rather than collapsing each run to a single global optimum; (ii) constructing an ensemble-level peak catalogue from the run-level goodness surfaces using objective prominence, width, separation, and reproducibility criteria; and (iii) explicitly allowing unresolved outcomes where the ensemble surface is flat, monotonic, boundary-dominated, or insufficiently supported. The resulting workflow retains the proven ability of the CDC test to identify primary Pb-loss events while providing a conservative basis for reporting more than one disturbance age only where the Monte Carlo ensemble preserves reproducible multi-modality. We validate these changes using synthetic datasets and show that the updated workflow can recover multi-episode structure that is obscured by single-solution summaries \citep{Mathieson2025a,Mathieson2025b,Reimink2016,Reimink2025}.

% -------------------------------------------------------------------------
% Replace lines 198--209 (CDC enhancements)
% -------------------------------------------------------------------------
\subsection{CDC enhancements}\label{sec:algorithmic-ext}
The most recent implementation of the CDC workflow repeated the concordant--discordant comparison hundreds of times under Monte Carlo resampling but, nevertheless, each realisation collapsed to a single Pb-loss age \citep{Mathieson2025a}. The revised pipeline introduced here lifts this limitation through three successive enhancements: (i) retaining local minima in each run that pass prominence and width gates, rather than only the global minimum (Sect.~\ref{sec:peakdetection}); (ii) merging these run-level peaks into an ensemble catalogue using a range-agnostic picker that operates in grid units, imposes explicit shoulder-merge, edge-handling, and reproducibility rules, and reports empirical confidence intervals for retained peaks (Sect.~\ref{sec:peakcatalogue}); and (iii) treating weakly structured cases conservatively by allowing unresolved outcomes when the ensemble surface is effectively flat, monotonic, boundary-dominated, or fails the support filters (Sect.~\ref{sec:disc-contrast}). Default thresholds are specified as fractions of the grid length or of the trimmed ensemble dynamic range so that results remain invariant to the user’s age window. These enhancements are described in detail below.

% -------------------------------------------------------------------------
% Delete lines 211--278 entirely
% -------------------------------------------------------------------------

% -------------------------------------------------------------------------
% Replace lines 282--296 (opening of Per-run peak detection)
% -------------------------------------------------------------------------
We use $R$ Monte Carlo realisations and a regular age grid of $N$ nodes. Analytical uncertainties are propagated by Monte Carlo resampling and the K--S dissimilarity is computed on the regular age grid $\{t_i\}_{i=1}^N$. For each run $n$ we evaluate the raw and penalised CDC dissimilarities,
\begin{equation}
  D_n(t_i), \qquad D^{*}_n(t_i),
\end{equation}
and define the corresponding goodness curves
\begin{equation}
  S_{\mathrm{raw},n}(t_i)=1-D_n(t_i),\qquad
  S_{\mathrm{pen},n}(t_i)=1-D^{*}_n(t_i).
\end{equation}
Unless stated otherwise, peak detection and ensemble voting are performed on the penalised curve $S_{\mathrm{pen},n}(t)$.

% -------------------------------------------------------------------------
% Replace line 989 sentence (runtime)
% -------------------------------------------------------------------------
On the benchmark grid (\(n_{\mathrm{grid}}=200\)) with \(R{=}200\) Monte Carlo replicates, the updated CDC workflow (including per-run peak detection and the ensemble catalogue) completed typical Case~1 to~4 datasets in a median of 71\,\unit{s} end-to-end (overall range 52 to 81\,\unit{s} across tiers). By tier, end-to-end medians were: Tier~A 71\,\unit{s} (61 to 80\,\unit{s}), Tier~B 70\,\unit{s} (58 to 81\,\unit{s}), and Tier~C 71\,\unit{s} (52 to 79\,\unit{s}) (Table~\ref{tab:cdc_runtime}).

% -------------------------------------------------------------------------
% Replace line 1095 paragraph
% -------------------------------------------------------------------------
The new workflow differs in three ways. First, candidate modes are detected on the robust median ensemble curve $\tilde S(t)$ using range-normalised thresholds, reducing sensitivity to the modelling window. Second, each catalogue entry carries a support value, defined as the fraction of Monte Carlo runs that cast a vote for that apex within its voting window, thereby providing an internal reproducibility measure alongside ages and intervals. Third, unresolved outcomes are treated explicitly: if the ensemble surface is effectively flat, monotonic, boundary-dominated, or insufficiently supported, no ensemble peak is reported. This keeps the method conservative when the data do not justify discrete disturbance ages.

% -------------------------------------------------------------------------
% Replace line 1109 sentence
% -------------------------------------------------------------------------
These tier-wise patterns reflect differences in how CDC and DD summarise the same Pb-loss age profiles. The roles of the concordant reference, the invalid-age penalty, and the peak-selection criteria are discussed in Sect.~\ref{sec:disc-contrast}.

% -------------------------------------------------------------------------
% Replace lines 1169--1178 subsection
% -------------------------------------------------------------------------
\subsubsection{Conservative abstention and boundary handling}
No ensemble peak is reported when the median goodness surface is effectively flat or monotonic across the tested window, when run-level optima collapse overwhelmingly onto a modelling-window boundary, or when candidate modes fail the ensemble support and interval filters. These rules prevent broad shoulders, edge artefacts, and discretisation effects from being promoted to disturbance ages, and ensure that weakly structured samples remain explicitly unresolved rather than being forced into a nominal Pb-loss estimate.

% -------------------------------------------------------------------------
% Replace line 1226 paragraph
% -------------------------------------------------------------------------
A failure mode arises when two Pb-loss episodes share a common upper intercept, but the younger episode splays off the older chord. At one candidate age, many analyses from the other episode reconstruct to impossible ages and incur the invalid-age penalty. This can collapse the ensemble surface toward a single intermediate maximum even when both episodes are real. We therefore regard the single-upper-intercept splay as a structural limitation of the present CDC method. Addressing this case likely requires a model that allows different discordant sub-arrays to be evaluated against different candidate episodes simultaneously, rather than penalising feasibility under a single shared trial lower intercept. This behaviour is consistent with prior synthetic overprinting experiments showing that multiple Pb-loss events cannot be separated when no analyses retain a single-stage record of the older episode \citep{Kirkland2020}.

% -------------------------------------------------------------------------
% Replace lines 1231--1232 conclusion opening
% -------------------------------------------------------------------------
We present a concordant--discordant comparison (CDC) workflow that extends a single-optimum K--S grid search to a multi-peak ensemble summary. The workflow preserves run-level peaks that pass explicit width and prominence criteria, merges these candidates into an ensemble catalogue using reproducibility voting and range-normalised peak metrics, and explicitly allows unresolved outcomes where the goodness surface does not support discrete disturbance ages. Each reported peak is accompanied by an empirical 95\,\% interval and a support value that summarises peak stability under Monte Carlo resampling.
